\section{Methods}
\label{sec:methods}

\subsection{Models and inputs}
The main model is a U-Net~\cite{ronneberger2015unet} with four levels, base width 32 and channel
dropout 0.1 (7.8\,M parameters), trained from scratch. A ResNet-34 U-Net~\cite{he2016resnet} with
ImageNet weights (segmentation-models-pytorch~\cite{iakubovskii2019smp}) is the second backbone.
Section~\ref{sec:cd} adds a change-detection U-Net and Section~\ref{sec:fm} a TerraMind foundation
model. Inputs are scaled with fixed, data-independent transforms so that every event is treated
identically: VV and VH backscatter are clipped to $[-30, 5]$\,dB, shifted by $+12.5$\,dB and divided by 7.5\,dB;
HAND $h$ (metres) enters as $\log(1+h)/2$; the DEM enters as relief relative to the tile median, divided
by 50\,m and clipped to $\pm 3$; slope is divided by 15\textdegree. Pixels with any missing input are
ignored in training and scoring.

\subsection{Physics losses and post-processing}
All models minimize binary cross-entropy plus Dice loss. Two optional terms encode HAND as a physical
constraint on the output. The \emph{monotonicity} loss penalizes a pixel for being predicted wetter than
a pixel that lies lower above drainage: for pixel pairs $(a,b)$ a fixed stride apart along rows and
columns with $h_a - h_b \ge 1$\,m, it adds $\max(0, p_a - p_b - m)$ for predicted probabilities $p$ and
margin $m$, averaged over qualifying pairs. The \emph{ceiling} loss penalizes flood predicted above
15\,m HAND. As a post-processing baseline, every run is also scored with predictions above 15\,m HAND
set to dry.

\subsection{Training}
We use AdamW~\cite{loshchilov2019adamw} (learning rate $10^{-3}$, weight decay $10^{-4}$) with a one-cycle
schedule~\cite{smith2019superconvergence}, bf16 mixed precision, random flips and $90^\circ$ rotations,
and select the checkpoint with the best validation IoU. Sen1Floods11 models train for 100 epochs on
$256^2$ crops (batch 16) and validate on the official validation split or, for leave-one-event-out and
all-event training, on 15\% of the training chips. GEOID-Flood models
train on $512^2$ crops of the $1024^2$ tiles (batch 8) for a fixed budget of 32,670 steps---30 epochs of
the full training set---whatever the training subset, so every subset size receives the same compute,
and validate on GEOID-Flood's validation split. Training is seeded; cuDNN kernels are not forced to be
deterministic.

\subsection{Evaluation}
The flood event is the unit of analysis. For each event we sum confusion counts over its tiles and compute
IoU of the water (or flood) class; \emph{event-mean} IoU averages these over events, weighting each event
equally. \emph{Pooled} IoU sums confusion counts over all test events, GEOID-Flood's convention, and so
weights events by their flooded area. The two can disagree when a few large events dominate.

Variants are compared in pairs that share everything but the tested factor, including the seed (or the
training draw, Section~\ref{sec:geoid}). For each event we average the paired difference over seeds or
draws and report the mean over events with a 95\% confidence interval (CI). For multi-seed results
(Sections~\ref{sec:transfer}--\ref{sec:failures}) the CI is, unless stated otherwise, a two-way
bootstrap~\cite{efron1993bootstrap} that resamples both seeds or draws and test events (4,000--10,000
resamples), which accounts for training variance as well as event variance; for pooled IoU the bootstrap
resamples the same two factors and recomputes the pooled counts. Protocol-matched pooled differences
reported next to absolute \iouf{} values (Sections~\ref{sec:transfer}, \ref{sec:fm}, \ref{sec:geoid}
and \ref{sec:cd}) are computed from confusion counts averaged over seeds or draws, so they can differ
slightly from the difference of the mean per-seed pooled IoUs beside them; all other pooled differences
are averages of per-seed differences. Average-precision CIs and the JRC comparisons in the Sen1Floods11
transfer resample test events only, as noted where they appear. Where indicated, we add a Wilcoxon signed-rank
test~\cite{wilcoxon1945} over events. Correlations between per-event gains and event descriptors use
Spearman's $\rho$ with Benjamini--Hochberg correction~\cite{benjamini1995fdr}.

\emph{Multiple comparisons.} The study makes many paired comparisons, and apart from the correlations the
CIs are not corrected for multiplicity. We therefore read a result as established only when its CI is
well clear of zero or when it was pre-registered; differences whose CI only just excludes zero (lower
bound within about 0.005 of it) are called suggestive, and the main conclusions do not rest on them. The
central transfer result ($p < 10^{-4}$ over 79 events) is far from this boundary.

\emph{Assumptions.} The bootstrap treats test events as exchangeable; events from the same country may not
be, so we also resample countries (Section~\ref{sec:transfer}). In the learning curve, draw $k$ is trained
with seed $k$, so draw and initialization variance are confounded, and draws of one size share events: the
CIs describe variation over draw--seed pairs, not over independent datasets.

\subsection{GEOID-Flood protocol}
GEOID-Flood scores single-image models by running them on the pre- and the post-event image and calling
flood what is water after but not before; permanent water counts as non-flood, and confusion counts are
pooled over the test set. We reproduce this protocol (using the latest pre-event Sentinel-1 image per
tile) and report \iouf{} under it alongside our own label modes. Our test set differs slightly from
GEOID-Flood's (EMSR388 and EMSR853 excluded), so comparisons with their numbers are indicative.

\subsection{Pre-registration}
Eight analyses were fixed in writing, with the analysis code committed to the repository, before their
outcomes were computed: the confirmation test of the learning curve (Section~\ref{sec:geoid}), the
relief gate and HAND dropout (Section~\ref{sec:failures}), the terrain ablation and TerraMind runs on the
Sen1Floods11 transfer (Sections~\ref{sec:ablation} and \ref{sec:fm}), the HAND perturbation tests
(Section~\ref{sec:mechanism}), and the stratified errors and country-cluster bootstrap
(Section~\ref{sec:transfer}). The plans and outcomes are listed in the repository's
\texttt{PREREGISTRATIONS.md}. Post hoc analyses are marked
exploratory.
